\section{Statistical model and main theorem}\label{sec:model}

\subsection{Persistent coins and the observation budget}
Let
\[
 \mu=\sum_{i=1}^n w_i\delta_{d_i},\qquad
 0<d_1<\cdots<d_n\le\tfrac12,\qquad
 w_i>0,\quad\sum_{i=1}^n w_i=1,
\]
be a known finite probability law. Equal gaps are merged before using
this representation. A parameter is a function $q$ on the support of
$\mu$. The complete latent pairs $(D_k,S_k)$ of distinct fresh coins
are iid, with $D_k\sim\mu$, $S_k\in\{-1,+1\}$, and
$\Pp_q(S_k=+1\mid D_k=d)=q(d)$. Each coin retains its pair on every
revisit. Conditional on all latent pairs, the observation tapes of
different coins are independent, and tape $k$ consists of iid
Bernoulli$(1/2+S_kD_k)$ observations. Neither $D_k$ nor $S_k$ is observed.

Fix a known rarity bound $0<\rho\le1/4$. The parameter class and target are
\begin{equation}\label{eq:parameter-class}
 \mathcal Q_\rho=\left\{q:\ 0\le q(d_i)\le1,\quad
       \theta(q):=\sum_{i=1}^n w_iq(d_i)\le\rho\right\}.
\end{equation}
This is a global constraint on the positive-sign fraction; it does not
impose $q(d_i)\le\rho$ at each gap.

An algorithm may create, revisit and interleave coins, and use an
independent random seed whose distribution does not depend on $q$.
Its choices and output are measurable functions of that seed, coin
identities and purchased observations. It must terminate and return
a real estimate $\widehat\theta$ almost surely under every
$q\in\mathcal Q_\rho$. Each observed flip costs one, including the
first observation of a fresh coin. Let $N$ be the total number of
observations. Coin creation, arithmetic and random-number generation
are not charged; creating an unobserved coin yields no information.
Zero observations are allowed, and no deterministic total cap is imposed.

For a known budget $B>0$, define the minimax root-mean-square error
\begin{equation}\label{eq:minimax-risk}
 R_{\mu,\rho}(B)=
 \inf_{\mathcal A:\ \sup_{q\in\mathcal Q_\rho}\E_qN\le B}
 \ \sup_{q\in\mathcal Q_\rho}
       \left\{\E_q\bigl(\widehat\theta-\theta(q)\bigr)^2\right\}^{1/2}.
\end{equation}
The budget is uniform over the entire parameter class. A procedure may
depend on $(\mu,\rho,B)$, but not on $q$ or any genuine coin's latent
state. Projecting its estimate onto $[0,\rho]$ cannot increase loss
or work. We abbreviate $R_{\mu,\rho}(B)$ to $R(B)$ when the inputs
are fixed.

\subsection{The closed quantile functional}
Remove exactly the mass $\rho$ with smallest gaps. In terms of the
weights, set
\begin{equation}\label{eq:quantile-split}
 \tau_i=\min\left\{w_i,\left(\rho-\sum_{\ell<i}w_\ell\right)_+\right\},
 \qquad \nu=\sum_i\tau_i\delta_{d_i},\qquad \kappa=\mu-\nu.
\end{equation}
Here $x_+=\max\{x,0\}$. Let $t$ be the smallest support point for
which $\mu((0,t])\ge\rho$. Then $\nu$ has mass $\rho$ and is supported
on $(0,t]$, while $\kappa$ has mass $1-\rho$ and is supported on
$[t,1/2]$. The atom at $t$ may be split fractionally.

For $j=0,1,2,\ldots$, define the dyadic bands and exact core costs by
\begin{equation}\label{eq:core-entropy}
 \begin{gathered}
 b_j=2^{-j-1},\qquad I_j=(b_j/2,b_j],\qquad
 a_j^*=\sum_{i:d_i\in I_j}\frac{w_i-\tau_i}{d_i^2},
 \qquad A^*=\sum_j a_j^*,\\
 \Gamma_\kappa=A^*+\sum_{j:a_j^*>0}a_j^*\log(A^*/a_j^*).
 \end{gathered}
\end{equation}
All logarithms are natural. There are finitely many positive costs,
and $A^*>0$. Thus $\Gamma_\kappa=A^*(1+\Ent(a^*/A^*))$ is the
inverse-square cost of the core multiplied by one plus the entropy
of its distribution across dyadic scales.

With $h=B^{-1/2}$, set
\begin{equation}\label{eq:closed-functional}
 T_h=\sum_i\tau_i\min\{1,h/d_i\},\qquad
 \Q(B)=\min\left\{\rho,\ T_h+h\sqrt{\rho\Gamma_\kappa}\right\}.
\end{equation}
The expression depends only on the known weights, gaps, rarity and
budget. No optimization over subsets or submeasures is needed to
evaluate it.

\subsection{The theorem and its quantifiers}
\begin{theorem}[Closed quantile characterization]\label{thm:direct-main}
For every integer $n\ge1$, every finite law $\mu$ described above,
every $0<\rho\le1/4$, and every $B>0$,
\begin{equation}\label{eq:main-bounds}
       10^{-20}\Q(B)\ \le\ R_{\mu,\rho}(B)\ \le\ 30001\Q(B).
\end{equation}
More precisely, for each $(\mu,\rho,B)$ there is one observable
algorithm satisfying, simultaneously for all $q\in\mathcal Q_\rho$,
\[
 \E_qN\le B,\qquad
 \left\{\E_q(\widehat\theta-\theta(q))^2\right\}^{1/2}
 \le30001\Q(B).
\]
Conversely, every algorithm with
$\sup_{q\in\mathcal Q_\rho}\E_qN\le B$ has worst-case RMS error
at least $10^{-20}\Q(B)$, including algorithms with arbitrary
adaptive interleaving, nonlinear outputs and unbounded stopping.
\end{theorem}

The constants are universal and are not optimized. The claim is a
comparison up to constants, not an exact formula for the minimax
coefficient. The iid assumption concerns the complete latent pairs,
not merely their magnitudes. The theorem concerns known finite laws
and expected observation budgets; other information structures and
budget conventions are not included.

\subsection{Notation and testing conventions}\label{sec:notation}
Write $s=\mathbf1_{\{S=+1\}}\in\{0,1\}$ for a binary sign label;
thus the conditional observation mean is $1/2+(2s-1)d$. A
single-unit procedure means a procedure on one genuine coin, with
its latent state retained throughout. A null prior is $q\equiv0$;
$\E_0$ denotes expectation at that prior. When the two signs are
compared at one specified gap, $\E_{-d}$ and $\E_{+d}$ refer to
Bernoulli means $1/2-d$ and $1/2+d$, respectively. Fair-coin
expectations are explicitly labelled $\E_{\rm fair}$.

For probability laws $P,Q$, dominated by a common measure, write
\[
 \Aff(P,Q)=\int\sqrt{dP\,dQ},\qquad
 \mathfrak h(P,Q)=1-\Aff(P,Q),\qquad
 \TV(P,Q)=\sup_E|P(E)-Q(E)|.
\]
In particular, $\mathfrak h$ denotes squared Hellinger distance in
the convention used here, whereas $h=B^{-1/2}$ is a budget scale.
We use $\TV\le\sqrt{2\mathfrak h}$, $\mathfrak h\le\KL/2$ and
Pinsker's inequality $\TV\le\sqrt{\KL/2}$, as well as KL data
processing and the finite-prefix chain rule. Unbounded stopping
is handled by finite caps and the stated limiting arguments.
For a probability vector $p$, $\Ent(p)=\sum_{p_j>0}p_j\log(1/p_j)$;
$\log_+x=\max\{0,\log x\}$, with $\log_+0=0$. The symbols
$\lesssim$, $\gtrsim$ and $\asymp$ suppress universal positive
constants only.

\begin{lemma}[Two-point RMS testing]\label{lem:two-point}
If two admissible parameters have target separation $\Delta>0$
and terminal transcript laws of total variation at most $1/2$,
every estimator has maximum RMS error at those parameters at
least $\Delta/4$.
\end{lemma}
\begin{proof}
Threshold the estimate at the midpoint of the two target values.
Each testing error incurs squared estimation loss at least
$\Delta^2/4$. The sum of the two testing errors is at least
$1-\TV\ge1/2$. Hence the sum of the two mean-square errors is
at least $\Delta^2/8$, and their maximum is at least $\Delta^2/16$.
\end{proof}
